\subsection{Overall Performance Evaluation}

The performance of RiskAdaptive was evaluated against conventional
machine learning baselines using the TON\_IoT network telemetry
dataset. Logistic Regression and Random Forest models were considered
as baseline intrusion detection approaches.

\begin{table}[H]
\caption{Performance Comparison with Baseline Models}
\centering
\begin{tabular}{lcccc}
\toprule
\textbf{Model} &
\textbf{Accuracy} &
\textbf{Precision} &
\textbf{Recall} &
\textbf{F1-score}\\
\midrule
Logistic Regression & 98.87\% & 99.30\% & 99.25\% & 99.28\%\\
Random Forest & 99.27\% & 99.94\% & 99.15\% & 99.54\%\\
RiskAdaptive & 99.27\% & 99.94\% & 99.15\% & 99.54\%\\
\bottomrule
\end{tabular}
\end{table}


The results demonstrate that RiskAdaptive maintains competitive
classification performance while extending conventional IDS capability
through risk-aware analysis. Unlike traditional classifiers that
produce only attack labels, the proposed framework additionally
generates security priorities based on estimated risk and severity.


\subsection{Attack-Specific Analysis}

To evaluate robustness across different threat categories, the proposed
framework was analyzed against multiple attack scenarios including
backdoor, ransomware, and man-in-the-middle attacks.

Backdoor attacks achieved high detection capability with elevated risk
scores, indicating effective identification of severe attack patterns.
Ransomware traffic was also successfully detected with high confidence,
demonstrating the capability of RiskAdaptive to prioritize impactful
security events.

Man-in-the-middle attacks showed comparatively lower detection
performance due to their stealth characteristics and traffic similarity
with legitimate communication patterns. This highlights the challenge
of detecting low-visibility attacks in IoT environments.


\subsection{Risk Prioritization Capability}

Beyond detection accuracy, RiskAdaptive provides an additional
decision-support layer by categorizing detected events according to
severity levels. The framework separates security events into HIGH,
MEDIUM, and LOW risk categories, enabling analysts to prioritize
critical incidents.

This capability addresses a limitation of conventional IDS systems,
where every detected attack is generally treated with equal priority.
By incorporating contextual risk scoring, RiskAdaptive improves the
practical usability of automated security monitoring.


\subsection{Discussion}

The experimental findings indicate that combining machine learning
based intrusion detection with agentic risk analysis provides a more
complete security workflow. While baseline models demonstrate strong
classification performance, they lack mechanisms for contextual
decision making.

RiskAdaptive extends existing IDS approaches by transforming detection
outputs into actionable security intelligence. The framework therefore
provides a bridge between automated attack detection and intelligent
incident response for Edge-IoT environments.
